\documentclass{article}
\usepackage[T1]{fontenc}
\usepackage{lmodern}
\usepackage{graphicx}
\usepackage{amsmath,amssymb}
\usepackage[a4paper,margin=2cm]{geometry}
\usepackage[absolute,overlay]{textpos}
\setlength{\TPHorizModule}{1mm}
\setlength{\TPVertModule}{1mm}
\usepackage[indonesian]{babel}
\usepackage{hyperref}
\setlength{\emergencystretch}{2em}
% unit: Halaman 3

\begin{document}

% === Halaman 3 (llm) ===

\begin{center}
\textbf{A PUBLIC KEY CRYPTOSYSTEM AND A SIGNATURE SCHEME BASED ON DISCRETE LOGARITHMS}\
\vspace{1em}
Taher ElGamal$^*$\
\vspace{1em}
Hewlett-Packard Labs\\
1501 Page Mill Rd.\\
Palo Alto CA 94301
\end{center}

\textbf{ABSTRACT}

A new signature scheme is proposed together with an implementation of the Diffie - Hellman key distribution scheme that achieves a public key cryptosystem. The security of both systems relies on the difficulty of computing discrete logarithms over finite fields.

\textbf{1. INTRODUCTION}

In 1976, Diffie and Hellman [3] introduced the concept of public key cryptography. Since then, several attempts have been made to find practical public key systems (see for example [6,7,9]) depending on the difficulty of solving some problems. For example, the RSA system [9] depends on the difficulty of factoring large integers. This paper presents systems that rely on the difficulty of computing logarithms over finite fields.

Section 2 shows a way to implement the public key distribution scheme introduced by Diffie and Hellman [3] to encrypt and decrypt messages. The security of this system is equivalent to that of the distribution scheme. Section 3 introduces a new digital signature

\end{document}
